\section{9. Creating and balancing encounters}

\subsection{The dungeon as a whole}

Difficulty comes from cumulative exposure to danger. A single encounter need not threaten the entire party. Repeated encounters, expended spell slots, damaged characters, consumed items and unreliable escape routes create pressure naturally. Because eleven ordinary rooms are dangerous, making every dangerous encounter an unavoidable lethal battle would be excessively punishing. Players must be able to recognize threats, retreat and make informed decisions.

\subsection{Recommended difficulty framework}

\begin{longtable}{p{0.3\linewidth}p{0.17\linewidth}p{0.45\linewidth}}

\rowcolor{tablehead} \textbf{Challenge} & \textbf{Suggested DC} & \textbf{Intended pressure} \\

\endfirsthead

\rowcolor{tablehead} \textbf{Challenge} & \textbf{Suggested DC} & \textbf{Intended pressure} \\

\endhead

Routine inspection & 10–12 & Rewards attention \\

Ordinary lock or trap & 13–15 & Moderate obstacle \\

Difficult mechanism & 16–18 & Requires skill or resources \\

Severe hazard & 18–20 & Significant danger \\

Exceptional challenge & 21+ & Rare, extraordinary obstacle \\

\end{longtable}

These are starting points, not mandatory statistics. Use encounter-building guidance appropriate to your D\&D edition and party. Do not rely exclusively on Challenge Rating: action economy, resistances, battlefield conditions and remaining resources can substantially affect difficulty.

\subsection{Easy fights}

Easy fights should generally be manageable without substantial resource expenditure. They create tension, reward tactical play and make collecting a chest slightly more difficult than walking across the room. Suitable opponents include skeletons, zombies, giant rats, spiders, animated weapons, imps and minor constructs. Use straightforward behavior and short encounters. Do not increase enemy numbers without considering their combined actions.

\subsection{Hard fights}

Hard fights represent meaningful threats: powerful enemies, coordinated groups, spellcasters or creatures supported by environmental hazards. Examples include necromancers with undead servants, elemental guardians, infernal creatures, troll packs and animated stone soldiers. These encounters contain no ordinary reward chests, so players have little incentive to fight unnecessarily. Their principal reward is surviving, discovering the room and potentially securing a useful route.

\subsection{Boss fights}

Bosses are particularly dangerous but provide a valuable reward. Give each boss a distinctive weakness or environmental mechanic that players can discover and exploit.

The weakness should provide a meaningful advantage, such as reducing defenses, interrupting a powerful ability or removing regeneration. Avoid making it an arbitrary mandatory solution. A capable party should still be able to defeat the boss through conventional combat. Give observable clues encouraging experimentation.

\subsection{Example: Clockwork Tyrant}

The characters enter a chamber dominated by an enormous mechanical guardian. Gears turn along the walls, and thick metal shafts connect machinery to its armored body. For a mid-level party, use a suitable construct stat block and add an environmental mechanism: the guardian gains protection while machinery operates; characters can investigate, attack or obstruct exposed gears; disabling machinery removes protection; the guardian attempts to prevent interference. Select exact statistics according to party strength. Fighters, spellcasters and technically skilled characters can all contribute.
